% preamble.tex -- ваше місце для ДОДАТКОВИХ пакетів і власних команд.
%
% kpithesis.cls навмисно не змінюють напряму (щоб оновлення шаблону
% в майбутньому не конфліктували з вашими правками) -- усе, що
% специфічне саме для ВАШОЇ роботи (додаткові пакети, власні
% математичні позначення, налаштування лістингів коду тощо), додавайте
% сюди. Цей файл підключається в main.tex одразу після
% \documentclass, тобто ще до \begin{document} -- як і будь-який
% "преамбульний" код.
%
% Порожній файл нічого не ламає -- усе нижче вимкнене коментарями.
% Розкоментуйте те, що вам потрібно, або дописуйте своє за тим самим
% зразком.

% --- Приклад: лістинги програмного коду (пакет listings) ---
% Якщо у вашій роботі є фрагменти коду -- розкоментуйте цей блок. Він
% налаштовує зовнішній вигляд лістингів (рамка, номери рядків, той
% самий моношрифт, що вже підключений шаблоном -- див.
% docs/ARCHITECTURE.md, розділ "Шрифти").
%
% \usepackage{listings}
% \lstset{
%   basicstyle       = \ttfamily\small,
%   numbers          = left,
%   numberstyle      = \tiny,
%   frame            = single,
%   breaklines       = true,
%   showstringspaces = false,
% }
%
% Використання в тексті розділу:
%   \begin{lstlisting}[language=Python, caption={Приклад функції}]
%   def hello():
%       print("Привіт, світ!")
%   \end{lstlisting}

% --- Приклад: власна математична команда ---
% \newcommand{\R}{\ensuremath{\mathbb{R}}}   % \R замість \mathbb{R}

% --- Приклад: додатковий пакет загального призначення ---
% \usepackage{tikz}   % для схем і діаграм
